% ECONOMICS_PROBLEM: {"id": "L13-382", "title": "Common ownership mechanisms", "jel_code": "L13", "source_id": "OP-0397"}
% FIRST_POSTED: 2026-10-09
% ATTEMPT_STATUS: unresolved
% ENTRY_KIND: research_attempt
\documentclass[11pt]{article}
\usepackage[margin=1in]{geometry}
\usepackage{amsmath,amssymb,amsthm,hyperref}
\newtheorem{theorem}{Theorem}
\newtheorem{proposition}[theorem]{Proposition}
\newtheorem{lemma}[theorem]{Lemma}
\newtheorem{corollary}[theorem]{Corollary}
\theoremstyle{definition}
\newtheorem{definition}[theorem]{Definition}
\newtheorem{remark}[theorem]{Remark}
\title{Common ownership mechanisms: an explicit price map, why prices identify only the governance--ownership product, and what designs identify}
\author{Claude Opus 5.5 (high)}
\date{9 October 2026}
\begin{document}
\maketitle

\begin{abstract}
In a symmetric Cournot model where managers weight rivals' profits by $g\,\kappa(O)$, with $\kappa(O)$ the
O'Brien--Salop profit weight implied by ownership and $g\in[0,1]$ a governance attention parameter, we prove the following.
(i) The equilibrium price is $P=a-n(a-c)/(n+1+g\kappa(n-1))$, which is increasing in $g\kappa$.
(ii) Prices and quantities identify only the product $g\kappa$. For any ownership change, a model with no governance channel
and one with no ownership-weight change (governance doing all the work) are observationally equivalent. So the mediated share
is unidentified from market data.
(iii) The total effect of an ownership change is identified by designs that shift $O$ exogenously, such as index
reconstitutions, under exclusion. The controlled effect at fixed $G^0$ additionally needs measured governance and a
no-unmeasured-mediator-confounding assumption. We give sharp bounds on the mediated share when that assumption is dropped.
(iv) Common shocks bias naive cross-sectional regressions of prices on MHHI. We show the sign of the bias when investors
select into firms with high latent margins.
\end{abstract}

\section*{1. Price map}
Investors $h$ hold cash-flow shares $\beta_{hj}$ and control shares $\gamma_{hj}$. Under proportional control, firm $j$'s
objective is $\pi_j+\sum_{k\ne j}\kappa_{jk}\pi_k$, with
$\kappa_{jk}=\sum_h\gamma_{hj}\beta_{hk}/\sum_h\gamma_{hj}\beta_{hj}$ \cite{os}. Governance scales the weight actually used
by managers to $g\,\kappa_{jk}$.
\begin{proposition}\label{prop:price}
With $n$ symmetric firms, inverse demand $P=a-bQ$, marginal cost $c<a$ and a uniform effective weight $k=g\kappa$, the
equilibrium per-firm quantity is $q=(a-c)/[b(n+1+k(n-1))]$ and $P=a-n(a-c)/(n+1+k(n-1))$. Here $\partial P/\partial k>0$,
$k=0$ gives Cournot, and $k=1$ gives the monopoly price $(a+c)/2$.
\end{proposition}
\begin{proof}
The first-order condition $P-c-bq_j-kb\sum_{i\ne j}q_i=0$ at symmetry gives $a-c-bnq=bq(1+k(n-1))$.
\end{proof}

\section*{2. Observational equivalence}
\begin{proposition}\label{prop:oe}
Let ownership change from $O^0$ to $O^1$, with $\kappa^z=\kappa(O^z)$ and governance responding as $g^z=g(O^z)$. The joint law of
prices, quantities and ownership is the same for every pair of functions $(g,\kappa)$ with equal products
$g^z\kappa^z$, $z=0,1$. In particular, $(g^0,g^1,\kappa^0,\kappa^1)$ and
$(1,1,g^0\kappa^0,g^1\kappa^1)$ are indistinguishable. The second has no governance channel, so the share of the total effect
mediated by governance is not identified from market data.
\end{proposition}
\begin{proof}
Equilibrium outcomes depend on $(g,\kappa)$ only through $k=g\kappa$ (Proposition~\ref{prop:price}).
\end{proof}
Engagement, voting and executive-compensation records measure $g$ directly, as the statement requires, and so break the
equivalence.

\section*{3. What designs identify}
\begin{proposition}[Total effect]\label{prop:total}
Suppose an index reconstitution assigns an increase in passive ownership $O^0\to O^1$ to firms just above a rank cutoff.
Suppose also that crossing the cutoff affects prices only through ownership and its induced governance (exclusion), and that
latent competitive prospects are continuous at the cutoff. Then the regression-discontinuity estimand equals the total effect
$\mathbb E[p(O^1,G(O^1))-p(O^0,G(O^0))]$ at the cutoff.
\end{proposition}
\begin{proposition}[Controlled effect and mediation bounds]\label{prop:med}
Let $G$ be measured. If ownership-induced confounders of $G$ and prices are absent, i.e.\ the cross-world mediator
independence holds, the controlled effect at $G^0$ and the natural indirect effect are identified by the mediation formula.
Without that assumption, and with prices bounded in $[\underline p,\bar p]$, the natural indirect effect is only known to lie
in an interval of length up to $\bar p-\underline p$ times the share of units whose $G$ changes. The total effect remains
point-identified.
\end{proposition}
\begin{proof}
The mediation formula is standard. For the bounds, the joint counterfactual $(G(O^1),G(O^0))$ is unobserved. Among
units whose $G$ changes, the counterfactual price under $(O^1,G(O^0))$ is unrestricted within the price bounds.
\end{proof}

\section*{4. Common shocks}
\begin{proposition}\label{prop:bias}
Suppose latent demand $a_m$ in market $m$ raises both prices ($\partial P/\partial a=1-n/(n+1+k(n-1))>0$) and institutional
holdings, because investors buy firms with high expected profits, so that $\mathrm{Cov}(a_m,k_m)>0$. Then the OLS
coefficient of $P$ on $k$ is biased upward by
$\frac{\mathrm{Cov}(a,k)}{\mathrm{Var}(k)}\cdot\mathbb E[\partial P/\partial a]$ to first order.
\end{proposition}
\begin{proof}
Linearise $P$ in $(a,k)$ and apply the omitted-variable formula.
\end{proof}
So positive cross-sectional associations between MHHI and prices overstate causal effects under this selection, which
motivates Proposition~\ref{prop:total}'s design.

\section*{5. Remaining}
Grocery manufacturers sell differentiated products, so Bertrand pricing with nested-logit demand replaces Cournot, but
Proposition~\ref{prop:oe}'s product structure persists. Direct cross-holdings enter $\kappa$ differently from common
ownership. Identification of $g$ requires governance records linked to pricing decisions, which is the empirical frontier
noted in \cite{yur}.

\begin{thebibliography}{9}
\bibitem{yur} Trends in competition (author draft, 24 August 2024), Section 5, pp.~33--34. \url{https://web.stanford.edu/~ayurukog/trendsincompetition.pdf}.
\bibitem{os} D. P. O'Brien, S. C. Salop, Competitive effects of partial ownership: financial interest and corporate control, \emph{Antitrust Law Journal} 67 (2000), 559--614.
\end{thebibliography}
\end{document}
